%% ============================================================
%% Rapport technique — Quantum Multiple Kernel Learning (QMKL)
%% Projet-QMKL-Finance (v1) + Projet-QMKL-v2
%% ============================================================

\documentclass[12pt,a4paper]{article}

% ----- Packages -----
\usepackage[utf8]{inputenc}
\usepackage[T1]{fontenc}
\usepackage[french]{babel}
\usepackage{amsmath,amssymb,amsthm}
\usepackage{mathtools}
\usepackage{bm}
\usepackage{geometry}
\usepackage{graphicx}
\usepackage{subcaption}
\usepackage{booktabs}
\usepackage{longtable}
\usepackage{array}
\usepackage{multirow}
\usepackage{xcolor}
\usepackage{hyperref}
\usepackage{listings}
\usepackage{algorithm}
\usepackage{algpseudocode}
\usepackage{tikz}
\usetikzlibrary{babel}
\usepackage{enumitem}
\usepackage{float}
\usepackage{fancyhdr}
\usepackage{titlesec}
\usepackage{tcolorbox}
\geometry{margin=2.5cm, headheight=15pt}

% ----- Colors -----
\definecolor{qblue}{RGB}{52,152,219}
\definecolor{qred}{RGB}{231,76,60}
\definecolor{qgreen}{RGB}{46,204,113}
\definecolor{qorange}{RGB}{243,156,18}
\definecolor{qpurple}{RGB}{155,89,182}
\definecolor{qgray}{RGB}{127,140,141}
\definecolor{codebg}{RGB}{248,248,248}
\definecolor{codeframe}{RGB}{200,200,200}

% ----- Code listings -----
\lstdefinestyle{pythonstyle}{
  backgroundcolor=\color{codebg},
  commentstyle=\color{qgray}\itshape,
  keywordstyle=\color{qblue}\bfseries,
  numberstyle=\tiny\color{qgray},
  stringstyle=\color{qgreen},
  basicstyle=\ttfamily\footnotesize,
  breakatwhitespace=false,
  breaklines=true,
  captionpos=b,
  keepspaces=true,
  numbers=left,
  numbersep=5pt,
  showspaces=false,
  showstringspaces=false,
  showtabs=false,
  tabsize=2,
  frame=single,
  rulecolor=\color{codeframe},
  language=Python
}
\lstset{style=pythonstyle}

% ----- Headers/footers -----
\pagestyle{fancy}
\fancyhf{}
\rhead{\textit{Quantum Multiple Kernel Learning}}
\lhead{\textit{Rapport technique}}
\rfoot{\thepage}
\lfoot{\textit{Projet-QMKL — 2025}}

% ----- Theorem environments -----
\newtheorem{definition}{Définition}[section]
\newtheorem{proposition}{Proposition}[section]
\newtheorem{remark}{Remarque}[section]

% ----- Custom boxes -----
\tcbuselibrary{skins,breakable}
\newtcolorbox{resultbox}[1][]{
  enhanced,
  title={#1},
  colback=qblue!5,
  colframe=qblue!50!black,
  fonttitle=\bfseries,
  breakable
}
\newtcolorbox{warningbox}[1][]{
  enhanced,
  title={#1},
  colback=qred!5,
  colframe=qred!50!black,
  fonttitle=\bfseries,
  breakable
}
\newtcolorbox{techbox}[1][]{
  enhanced,
  title={#1},
  colback=qgray!10,
  colframe=qgray!50!black,
  fonttitle=\bfseries,
  breakable
}

% ----- Hyperref -----
\hypersetup{
  colorlinks=true,
  linkcolor=qblue,
  citecolor=qgreen,
  urlcolor=qpurple,
  pdftitle={Rapport QMKL},
  pdfauthor={Projet-QMKL}
}

% ================================================================
\begin{document}

\begin{titlepage}
\centering
\vspace*{2cm}

{\Huge\bfseries Quantum Multiple Kernel Learning\par}
\vspace{0.5cm}
{\LARGE\bfseries pour la Classification Financière\par}
\vspace{1.5cm}

{\large\itshape Rapport technique complet\par}
{\large\itshape Projet-QMKL-Finance (v1) + Projet-QMKL-v2\par}

\vspace{2cm}
\rule{\textwidth}{1pt}
\vspace{0.5cm}
\begin{tabular}{ll}
\textbf{Cadre} & Projet académique — Informatique Quantique Appliquée \\
\textbf{Technologies} & Qiskit 1.x, scikit-learn, IBM Quantum (Open Plan) \\
\textbf{Matériel testé} & IBM Torino (133 qubits, Heron r2) \\
\textbf{Date} & Avril 2025 \\
\end{tabular}
\vspace{0.5cm}
\rule{\textwidth}{1pt}

\vspace{2cm}

\begin{abstract}
\noindent Ce rapport documente deux implémentations successives de Quantum Multiple Kernel Learning (QMKL)
pour la classification de données financières. Le \textbf{projet v1} (Projet-QMKL-Finance) réalise une
étude empirique exhaustive de 12 noyaux quantiques et 5 stratégies MKL sur trois jeux de données
standards (German Credit, Bank Marketing, Breast Cancer). Le \textbf{projet v2} (Projet-QMKL-v2)
étend l'approche avec une formulation QUBO pour l'assignation optimale de features aux noyaux,
testée sur données macroéconomiques FRED et sur matériel IBM Torino.

Les résultats empiriques sont mitigés : QMKL n'améliore pas significativement les classifieurs
classiques (RBF-SVM, Random Forest) dans nos conditions expérimentales. Ce rapport documente
honnêtement ce qui fonctionne, ce qui ne fonctionne pas, et \emph{pourquoi} — les résultats négatifs
étant eux-mêmes une contribution scientifique valide.
\end{abstract}

\vfill
\end{titlepage}

\tableofcontents
\newpage

% ================================================================
\section{Introduction et Contexte}
\label{sec:intro}

\subsection{Motivation}

Le Machine Learning Quantique (QML) promet un avantage algorithmique grâce à la superposition et
l'intrication quantiques. Parmi les approches QML, les \emph{noyaux quantiques} (quantum kernels)
présentent un attrait particulier : ils restent dans le paradigme SVM bien établi tout en projetant
les données dans des espaces de Hilbert exponentiellement grands, inaccessibles classiquement.

Les données financières semblent idéales pour ces approches :
\begin{itemize}
  \item Haute dimensionnalité (20–30 indicateurs macroéconomiques)
  \item Corrélations complexes non linéaires entre variables
  \item Faible rapport signal/bruit (crises rares, régimes changeants)
  \item Problèmes de classification binaire bien définis (récession / expansion, défaut / non-défaut)
\end{itemize}

\subsection{Objectifs des deux projets}

\paragraph{V1 — Projet-QMKL-Finance.}
Évaluer empiriquement si la combinaison de multiples noyaux quantiques (MKL) améliore la performance
de classification par rapport à (1) un noyau quantique unique et (2) des classifieurs classiques
de référence.

\paragraph{V2 — Projet-QMKL-v2.}
Si les barren plateaus limitent chaque noyau à $Q \leq 6$ qubits, comment couvrir optimalement
un espace de $d = 20$--30 features ? La réponse proposée est une formulation QUBO (Quadratic
Unconstrained Binary Optimization) pour l'assignation features $\to$ noyaux. Ce projet pose également
la question fondamentale : \emph{quand} l'avantage quantique se manifeste-t-il, et pour quel régime
de $Q$ ?

\subsection{Question de recherche centrale}

\begin{tcolorbox}[colback=qblue!10,colframe=qblue!60!black]
\textbf{Question centrale (v2)} : Comment assigner optimalement $d$ features à $M$ noyaux quantiques
via QUBO, et à partir de quel $Q$ (nombre de qubits par noyau) la simulation classique devient-elle
infaisable, justifiant l'usage d'un ordinateur quantique réel ?
\end{tcolorbox}

% ================================================================
\section{Fondements Théoriques}
\label{sec:theory}

\subsection{Noyaux Quantiques}

\subsubsection{Feature Map et Espace de Hilbert}

Un \emph{feature map quantique} encode un vecteur de données $\bm{x} \in \mathbb{R}^Q$ en un état
quantique $|\psi(\bm{x})\rangle \in \mathcal{H} = (\mathbb{C}^2)^{\otimes Q}$, un espace de
Hilbert de dimension $2^Q$ :
\begin{equation}
  U_\phi(\bm{x}) : |\bm{0}\rangle \mapsto |\psi(\bm{x})\rangle = U_\phi(\bm{x})|\bm{0}\rangle
\end{equation}

\subsubsection{Noyau de Fidélité}

Le noyau de fidélité mesure le recouvrement entre deux états quantiques :
\begin{equation}
  K_{\text{fid}}(\bm{x}, \bm{x}') = |\langle\psi(\bm{x}')|\psi(\bm{x})\rangle|^2
  = |\langle\bm{0}|U_\phi^\dagger(\bm{x}')U_\phi(\bm{x})|\bm{0}\rangle|^2
  \label{eq:fidelity}
\end{equation}
Ce noyau est valide au sens de Mercer (symétrique, semi-défini positif) et peut être estimé sur
un circuit quantique via la mesure du bit $|\bm{0}\rangle$ après $U_\phi^\dagger(\bm{x}')U_\phi(\bm{x})$.

\subsubsection{Noyau Projeté}

Le \emph{Projected Quantum Kernel} (PQK) réduit d'abord l'état à ses matrices densité
mono-qubit $\rho_k(\bm{x}) = \text{Tr}_{\neq k}|\psi(\bm{x})\rangle\langle\psi(\bm{x})|$,
puis applique un noyau RBF sur les distances de Frobenius :
\begin{equation}
  K_{\text{proj}}(\bm{x}, \bm{x}') = \exp\!\left(-\gamma \sum_{k=0}^{Q-1}
  \|\rho_k(\bm{x}) - \rho_k(\bm{x}')\|_F^2\right)
  \label{eq:projected}
\end{equation}

\subsection{Familles de Feature Maps Utilisées}

Deux familles principales ont été implémentées :

\subsubsection{ZFeatureMap (local, sans intrication)}
\begin{equation}
  U_Z(\bm{x}) = \bigotimes_{k=0}^{Q-1} R_Z(\alpha x_k)
  \label{eq:zmap}
\end{equation}
Le noyau correspondant admet une forme analytique fermée :
\begin{equation}
  K_Z(\bm{x}, \bm{x}') = \prod_{k=0}^{Q-1} \cos^2\!\left(\alpha(x_k - x'_k)\right)
  \label{eq:kz_analytical}
\end{equation}

\subsubsection{ZZFeatureMap (avec intrication)}
\begin{align}
  U_{ZZ}(\bm{x}) &= \left(\prod_{k=0}^{Q-2} \text{CX}_{k,k+1} \cdot R_Z\!\left(\alpha(x_k - x_{k+1})^2\right) \cdot \text{CX}_{k,k+1}\right) \cdot U_Z(\bm{x})
  \label{eq:zzmap}
\end{align}
Forme analytique (supposant pas d'intrication supplémentaire) :
\begin{equation}
  K_{ZZ}(\bm{x}, \bm{x}') = K_Z(\bm{x}, \bm{x}') \cdot
  \prod_{k<l} \cos^2\!\left(\alpha(x_k x_l - x'_k x'_l)\right)
  \label{eq:kzz_analytical}
\end{equation}

Ces formes analytiques ont été implémentées dans \texttt{src/kernels/analytical.py} (v2) pour
éviter la dépendance Qiskit dans les diagnostics, permettant des évaluations rapides sur CPU.

\subsection{Multiple Kernel Learning (MKL)}

\subsubsection{Combinaison linéaire}

Étant donnés $M$ noyaux $\{K_m\}_{m=1}^M$, la combinaison MKL est :
\begin{equation}
  K_{\text{comb}}(\bm{x}, \bm{x}') = \sum_{m=1}^M w_m K_m(\bm{x}, \bm{x}'),
  \quad w_m \geq 0, \quad \sum_{m=1}^M w_m = 1
  \label{eq:mkl_combination}
\end{equation}

\subsubsection{Alignement centré}

Le noyau idéal est le noyau cible $K_y[i,j] = \mathbb{1}[y_i = y_j]$.
L'alignement de noyau centré (Cortes \emph{et al.}, 2012) est :
\begin{equation}
  A(K_i, K_j) = \frac{\langle K_i^c, K_j^c \rangle_F}{\|K_i^c\|_F \cdot \|K_j^c\|_F}
\end{equation}
où $K^c = \left(I - \frac{\bm{1}\bm{1}^T}{m}\right) K \left(I - \frac{\bm{1}\bm{1}^T}{m}\right)$
est le noyau centré.

Le problème d'optimisation des poids est :
\begin{equation}
  \min_{\bm{v} \geq 0} \; \bm{v}^T M \bm{v} - 2\bm{v}^T \bm{a}
  \quad \text{où} \quad M_{ij} = \langle K_i^c, K_j^c\rangle_F, \quad a_i = \langle K_i^c, K_y^c\rangle_F
  \label{eq:centered_alignment}
\end{equation}

\subsubsection{Alignement SDP}

La formulation SDP (Lanckriet \emph{et al.}, 2004) :
\begin{equation}
  \max_{\bm{w}} \; \bm{w}^T \bm{q} \quad \text{s.t.} \quad \bm{w}^T S \bm{w} \leq 1, \; \bm{w} \geq 0
  \quad \text{où} \quad q_i = \langle K_i, K_y\rangle_F, \; S_{ij} = \langle K_i, K_j\rangle_F
  \label{eq:sdp_alignment}
\end{equation}
résolu via CVXPY avec cascade de solveurs : CLARABEL $\to$ ECOS $\to$ SCS.

\subsection{Barren Plateaus}
\label{sec:barren_plateaus}

\subsubsection{Phénomène}

Pour un circuit quantique aléatoire de $Q$ qubits, la variance du gradient d'un paramètre $\theta$
décroît exponentiellement :
\begin{equation}
  \text{Var}\!\left[\frac{\partial \mathcal{L}}{\partial \theta}\right] \sim O\!\left(2^{-Q}\right)
  \label{eq:barren_plateau}
\end{equation}
Cette concentration exponentielle rend l'optimisation par descente de gradient impossible pour
$Q \gtrsim 15$.

\subsubsection{Conséquence sur les noyaux}

Pour les noyaux de fidélité à grande profondeur, le phénomène analogue est la \emph{concentration
du noyau} : tous les $K(\bm{x}, \bm{x}')$ convergent vers une même constante (indépendamment de
$\bm{x}, \bm{x}'$), effaçant l'information discriminante.

\subsubsection{Mesure empirique}

Deux métriques ont été implémentées dans \texttt{src/kernels/diagnostics.py} :
\begin{itemize}
  \item \textbf{Variance du gradient} (finite-difference) :
  \begin{equation}
    \text{GradVar}(K, \alpha) = \text{Var}_{i \neq j}\!\left[\frac{K(X,X;\alpha+\epsilon) - K(X,X;\alpha-\epsilon)}{2\epsilon}\right]
  \end{equation}
  \item \textbf{Écart-type des entrées hors-diagonale} :
  \begin{equation}
    \sigma_{\text{off}}(K) = \text{Std}_{i \neq j}[K_{ij}]
  \end{equation}
\end{itemize}

\subsection{Formulation QUBO pour l'Assignation de Features}
\label{sec:qubo}

\subsubsection{Problème d'assignation}

\begin{definition}[Assignation Feature $\to$ Noyau]
Étant donnés $d$ features et $M$ noyaux quantiques de $Q$ qubits chacun, définir les variables
binaires $x_{k,m} \in \{0,1\}$ indiquant si la feature $k$ est assignée au noyau $m$.
L'assignation optimale satisfait exactement $Q$ features par noyau.
\end{definition}

\subsubsection{Formulation C (recommandée)}

\begin{equation}
\boxed{
\min_{\bm{x} \in \{0,1\}^{dM}}
  \underbrace{-\sum_{k,m} a_{k,m} \, x_{k,m}}_{\text{qualité}}
  + \underbrace{\lambda \sum_k \sum_{m < m'} x_{k,m} \, x_{k,m'}}_{\text{diversité}}
  + \underbrace{\mu_1 \sum_m \!\left(\sum_k x_{k,m} - Q\right)^{\!2}}_{\text{contrainte taille}}
}
\label{eq:qubo_C}
\end{equation}

\paragraph{Alignements marginaux $a_{k,m}$.}
Pour chaque paire (feature $k$, noyau $m$), l'alignement marginal mesure la qualité discriminante
du noyau $m$ en utilisant seulement la feature $k$ (avec $Q-1$ features de remplissage) :
\begin{equation}
  a_{k,m} = A(K_m^{(k)}, K_y^c)
\end{equation}
où $K_m^{(k)}$ est la matrice noyau du noyau $m$ construite uniquement avec la feature $k$.

\paragraph{Matrice QUBO.}
En développant \eqref{eq:qubo_C}, on obtient la représentation matricielle $E(\bm{x}) = \bm{x}^T Q \bm{x}$ :
\begin{align}
  Q_{ii} &= -a_{k,m} + \mu_1(1 - 2Q) \quad (i = k \cdot M + m) \\
  Q_{ij} &= 2\mu_1 \quad (i = k_1 M + m,\; j = k_2 M + m,\; k_1 < k_2) \\
  Q_{ij} &= \lambda \quad (i = k M + m_1,\; j = k M + m_2,\; m_1 < m_2)
\end{align}
$Q$ est triangulaire supérieure de taille $(dM) \times (dM)$.

\subsubsection{Solveurs QUBO}

Quatre solveurs ont été implémentés dans \texttt{src/qubo/solvers.py} (Tableau~\ref{tab:solvers}).

\begin{table}[H]
\centering
\caption{Comparaison des solveurs QUBO}
\label{tab:solvers}
\begin{tabular}{@{}lcccc@{}}
\toprule
\textbf{Solveur} & \textbf{Complexité} & \textbf{Optimalité} & \textbf{Faisable \#vars} & \textbf{Notes} \\
\midrule
Brute Force & $O(2^n)$ & Globale & $n \leq 24$ & Référence exacte \\
Greedy & $O(d^2 M)$ & Locale & Illimité & Très rapide \\
Recuit Simulé & $O(n_{\text{iter}})$ & Quasi-opt. & Illimité & $\approx$ optimal en pratique \\
QAOA ($p$ couches) & $O(n_{\text{circuit}})$ & Approx. & $\leq 50$ (sim.), illimité (HW) & Avantage quantique potentiel \\
\bottomrule
\end{tabular}
\end{table}

\subsubsection{QAOA — Quantum Approximate Optimization Algorithm}

QAOA résout des problèmes QUBO en convertissant la matrice $Q$ en Hamiltonien d'Ising :
\begin{equation}
  H_C = \sum_{i \leq j} Q_{ij} \, \frac{I - Z_i}{2} \cdot \frac{I - Z_j}{2}
  \label{eq:ising}
\end{equation}

L'ansatz QAOA à $p$ couches est :
\begin{equation}
  |\bm{\theta}\rangle = e^{-i\beta_p H_B} e^{-i\gamma_p H_C} \cdots e^{-i\beta_1 H_B} e^{-i\gamma_1 H_C} |+\rangle^{\otimes n}
  \label{eq:qaoa_ansatz}
\end{equation}
où $H_B = \sum_i X_i$ est l'Hamiltonien de mélange, $\bm{\theta} = (\bm{\beta}, \bm{\gamma})$
sont optimisés classiquement (COBYLA), et $n = dM$ variables binaires.

% ================================================================
\section{Projet V1 — QMKL-Finance}
\label{sec:v1}

\subsection{Architecture}

\begin{figure}[H]
\centering
\begin{tikzpicture}[
  box/.style={rectangle,draw,rounded corners,fill=qblue!15,minimum width=3cm,minimum height=1cm,align=center},
  arrow/.style={->,thick,qblue!70!black}
]
\node[box] (data) at (0,0) {Données brutes\\(OpenML, sklearn)};
\node[box] (preproc) at (4,0) {Prétraitement\\QuantumScaler + PCA};
\node[box] (kernels) at (8,0) {$M=12$ Noyaux\\Quantiques};
\node[box] (mkl) at (8,-2) {Combinaison MKL\\(Alignment/BO)};
\node[box] (svm) at (4,-2) {QSVM\\(noyau précalculé)};
\node[box] (eval) at (0,-2) {Évaluation\\AUC, F1, recall};

\draw[arrow] (data) -- (preproc);
\draw[arrow] (preproc) -- (kernels);
\draw[arrow] (kernels) -- (mkl);
\draw[arrow] (mkl) -- (svm);
\draw[arrow] (svm) -- (eval);
\end{tikzpicture}
\caption{Pipeline QMKL-Finance (v1)}
\label{fig:v1_pipeline}
\end{figure}

\subsection{Jeux de Données}

\begin{table}[H]
\centering
\caption{Jeux de données utilisés en v1}
\label{tab:v1_datasets}
\begin{tabular}{@{}llccccc@{}}
\toprule
\textbf{Dataset} & \textbf{Domaine} & $N$ & $d$ & \textbf{Classes} & \textbf{Équilibre} & \textbf{Source} \\
\midrule
German Credit & Finance & 1000 & 20 & 2 & 70/30\% & UCI/OpenML \#31 \\
Bank Marketing & Finance & 5000+ & 16 & 2 & 88/12\% & UCI/OpenML \#1461 \\
Breast Cancer & Médecine & 569 & 30 & 2 & 63/37\% & sklearn \\
\bottomrule
\end{tabular}
\end{table}

Les données sont prétraitées par \texttt{src/preprocessing/scaler.py} :
\begin{enumerate}
  \item Standardisation (zero-mean, unit-variance)
  \item PCA vers $Q$ composantes (réduction dimensionnelle)
  \item MinMax scaling vers $[0, 2\pi]$ (encodage quantique optimal)
\end{enumerate}

\subsection{Bibliothèque de 12 Noyaux Quantiques}

La bibliothèque \texttt{src/kernels/feature\_maps.py} (188 lignes) définit 12 configurations :

\begin{table}[H]
\centering
\caption{Bibliothèque de 12 noyaux quantiques (v1)}
\label{tab:v1_kernels}
\begin{tabular}{@{}clccc@{}}
\toprule
\textbf{ID} & \textbf{Famille Pauli} & \textbf{$\alpha$} & \textbf{Intrication} & \textbf{Profondeur} \\
\midrule
1–2 & Z (local) & 1.0 / 3.0 & Aucune & 1 \\
3–4 & ZZ (2-corps) & 1.0 / 4.0 & Linéaire & 1 \\
5–6 & XZ (croisé) & 0.5 / 2.5 & Linéaire, 2 reps & 2 \\
7–8 & YXX (3-corps) & 0.6 / 3.0 & Linéaire, 2 reps & 2 \\
9–10 & YZX (3-corps) & 0.6 / 3.0 & Linéaire, 2 reps & 2 \\
11–12 & Pauli (ZZ, 2 reps) & 0.6 / 2.5 & Linéaire, 2 reps & 2 \\
\bottomrule
\end{tabular}
\end{table}

La diversité des familles (Z local, ZZ 2-corps, XZ croisé, YXX/YZX 3-corps) vise à capturer
des interactions feature différentes dans l'espace de Hilbert.

\paragraph{Choix du paramètre $\alpha$.}
Le paramètre $\alpha$ contrôle la "résolution" du noyau. Une valeur faible ($\alpha \approx 0.5$)
produit un noyau très lisse (underfitting) ; une valeur élevée ($\alpha \approx 20$) produit un
noyau très concentré (approchant une identité). Les valeurs $\alpha \in \{0.5, 1.0, 2.5, 3.0, 4.0\}$
ont été sélectionnées empiriquement.

\subsection{Stratégies MKL}

Cinq stratégies de combinaison ont été comparées (Tableau~\ref{tab:v1_mkl_strategies}) :

\begin{table}[H]
\centering
\caption{Stratégies MKL implémentées}
\label{tab:v1_mkl_strategies}
\begin{tabular}{@{}llcc@{}}
\toprule
\textbf{Méthode} & \textbf{Principe} & \textbf{Complexité} & \textbf{Fichier src} \\
\midrule
Average & $w_m = 1/M$ (uniforme) & $O(M)$ & alignment.py \\
Centered & Minimisation quadratique \eqref{eq:centered_alignment} & $O(M^2)$ & alignment.py \\
SDP & Programme semi-défini \eqref{eq:sdp_alignment} & $O(M^3)$ & alignment.py \\
Projection & Greedy itératif sur résidus & $O(M^2)$ & alignment.py \\
Bayesian & BO sur poids + $C$ SVM & $O(n_{\text{calls}} \cdot M^2)$ & bayesian\_optimizer.py \\
\bottomrule
\end{tabular}
\end{table}

\subsubsection{Implémentation de l'Alignement Centré}

\begin{lstlisting}[caption=Alignement centré — src/mkl/alignment.py (extrait)]
def _center_kernel(K):
    """Centre le noyau dans l'espace des features."""
    n = K.shape[0]
    one = np.ones((n, n)) / n
    return K - one @ K - K @ one + one @ K @ one

def compute_centered_alignment_weights(kernels, K_target):
    """Minimisation QP : min v^T M v - 2 v^T a, v >= 0"""
    M_mat = np.array([[np.sum(Kc_i * Kc_j)
                       for Kc_j in kernels_c] for Kc_i in kernels_c])
    a_vec = np.array([np.sum(Kc_i * K_target_c) for Kc_i in kernels_c])

    # Regularisation: M <- M + eps * ||M||_F * I (stabilite numerique)
    M_mat += EPS_REG * np.linalg.norm(M_mat, 'fro') * np.eye(len(kernels))

    # Cascade de solveurs: CLARABEL -> ECOS -> SCS -> closed-form
    w = solve_qp(M_mat, a_vec)  # w >= 0
    return w / (w.sum() + 1e-12)  # Normalisation
\end{lstlisting}

\subsubsection{Optimisation Bayésienne}

La stratégie BO (\texttt{src/mkl/bayesian\_optimizer.py}) optimise conjointement :
\begin{itemize}
  \item Les poids $\bm{w} \in [0,1]^M$ (espace de Sobol)
  \item Le paramètre de régularisation $C \in [10^{-2}, 10^2]$ (log-uniforme)
\end{itemize}
via la minimisation de la validation croisée 5-fold, avec 50 appels au total (10 initiaux aléatoires,
40 par acquisition Gaussian Process). Cette méthode est environ 200$\times$ plus lente que
l'alignement centré pour un gain de +0.8 point AUC au mieux.

\subsection{Système de Cache de Noyaux}

\texttt{src/kernels/kernel\_matrix.py} (236 lignes) implémente :
\begin{itemize}
  \item \textbf{Cache disque} : hash MD5 des paramètres $\to$ \texttt{results/kernel\_cache/*.npy}
  \item \textbf{Calcul parallèle} : joblib ($n_{\text{jobs}} = -1$) pour $M$ noyaux simultanément
  \item \textbf{Régularisation} : $K \leftarrow K + \varepsilon I$ pour garantir $K \succ 0$
  \item \textbf{Diagnostics} : \texttt{kernel\_statistics()} détecte la concentration
\end{itemize}

\subsection{Résultats Empiriques}

\subsubsection{Configuration expérimentale}

\begin{itemize}
  \item $N = 200$ échantillons d'entraînement (sous-échantillonnage de German Credit/Bank Marketing)
  \item $Q = 6$ qubits par noyau
  \item $M = 12$ noyaux dans la bibliothèque
  \item $K$-fold $= 4$ splits, 20 tirages aléatoires (rapport mean $\pm$ std)
  \item Simulateur : Qiskit AER (statevector, shots = 8192)
\end{itemize}

\begin{table}[H]
\centering
\caption{Résultats comparatifs — AUC ROC (mean $\pm$ std)}
\label{tab:v1_results}
\begin{tabular}{@{}llccc@{}}
\toprule
\textbf{Dataset} & \textbf{Méthode} & \textbf{AUC} & \textbf{$\pm$ Std} & \textbf{$\Delta$ vs RBF} \\
\midrule
\multirow{7}{*}{German Credit}
  & Single K (Z, $\alpha=2$) & 0.712 & 0.067 & $-$12.3 \\
  & QMKL-Average            & 0.763 & 0.055 & $-$7.2 \\
  & QMKL-Centered           & 0.750 & 0.061 & $-$8.5 \\
  & QMKL-SDP                & 0.748 & 0.063 & $-$8.7 \\
  & QMKL-Projection         & 0.755 & 0.058 & $-$8.0 \\
  & QMKL-BO                 & 0.758 & 0.058 & $-$7.7 \\
  & \textbf{RBF-SVM}        & \textbf{0.835} & 0.042 & — \\
\midrule
\multirow{5}{*}{Bank Marketing}
  & QMKL-Average            & 0.741 & 0.151 & $-$7.6 \\
  & QMKL-BO                 & 0.774 & 0.135 & $-$4.3 \\
  & \textbf{RBF-SVM}        & \textbf{0.817} & 0.068 & — \\
  & Logistic Regression     & \textbf{0.867} & 0.057 & $+$5.0 \\
  & Random Forest           & 0.833 & 0.046 & $+$1.6 \\
\midrule
\multirow{4}{*}{Breast Cancer}
  & QMKL-Average            & 0.995 & 0.004 & $-$0.1 (n.s.) \\
  & QMKL-BO                 & 0.992 & 0.006 & $-$0.4 \\
  & \textbf{RBF-SVM}        & \textbf{0.996} & 0.004 & — \\
  & Random Forest           & 0.996 & 0.003 & $+$0.0 \\
\bottomrule
\end{tabular}
\end{table}

\begin{resultbox}[Résumé des résultats v1]
\begin{itemize}
  \item QMKL est systématiquement inférieur au RBF-SVM classique sur les deux jeux financiers
  \item L'écart est statistiquement significatif (test de Wilcoxon, $p < 0.001$) sur German Credit
  \item Sur Breast Cancer, les performances sont équivalentes ($p = 0.14$)
  \item La stratégie Average est souvent aussi bonne que les méthodes d'alignment sophistiquées
  \item La BO est $200\times$ plus lente pour $\leq +1$ point d'AUC
\end{itemize}
\end{resultbox}

\subsection{Analyses Avancées}

\subsubsection{Ablation sur Q (nombre de qubits)}

L'ablation du nombre de qubits ($Q \in \{2, 4, 6, 8, 10\}$) révèle un plateau de performance
dès $Q = 4$--6, puis une \emph{dégradation} au-delà (barren plateau se manifeste). Ce résultat
suggère que l'augmentation de $Q$ sans stratégie de mitigation est contre-productive.

\subsubsection{Interprétabilité : Valeurs de Shapley}

\texttt{src/mkl/shapley\_mkl.py} calcule les valeurs de Shapley des noyaux dans le QMKL,
quantifiant la contribution marginale de chaque noyau. Résultat notable : 3--4 noyaux dominent
et les 8--9 restants ont une contribution quasi-nulle, remettant en question l'utilité de
maintenir $M = 12$.

\subsubsection{MKL Riemannien}

\texttt{src/mkl/riemannian\_mkl.py} exploite la géométrie différentielle des matrices de noyaux
(variété riemannienne des matrices SDP). La distance géodésique entre noyaux est utilisée pour
pondérer la combinaison. Cette approche reste expérimentale et n'améliore pas significativement
les résultats sur nos données.

\subsection{Ce Qui N'a Pas Fonctionné en V1}

\begin{warningbox}[Résultats négatifs V1 — Enseignements]
\begin{enumerate}
  \item \textbf{Q élevé ($Q > 6$)} : La concentration exponentielle annule le gain théorique.
        La PCA pour réduire à $Q$ dimensions perd de l'information critique.

  \item \textbf{SDP} : Instabilité numérique fréquente (matrices SDP mal conditionnées).
        La cascade CLARABEL $\to$ ECOS $\to$ SCS $\to$ closed-form a été nécessaire.

  \item \textbf{Optimisation Bayésienne} : Le sur-ajustement sur la validation croisée et le coût
        computationnel ($50 \times$ évaluation noyau complète) ne justifient pas l'usage.

  \item \textbf{Noyau projeté} : Plus lent (calcule les matrices densité $\rho_k$) sans avantage AUC.

  \item \textbf{Bank Marketing} : Le déséquilibre de classes sévère (12\% positif) rend la
        comparaison par AUC trompeuse ; la régression logistique surpasse tout le monde.

  \item \textbf{MKL vs single kernel} : Sur Breast Cancer, un seul noyau $K_{ZZ}$ atteint
        AUC = 0.993, proche du QMKL. La complexité de combinaison ne se justifie pas.
\end{enumerate}
\end{warningbox}

% ================================================================
\section{Projet V2 — QMKL-v2 (Assignation QUBO)}
\label{sec:v2}

\subsection{Motivation : Dépasser les Barren Plateaus via l'Assignation}

Le projet v2 part d'un constat fondamental : avec $d = 19$--30 features et $Q \leq 6$ qubits
par noyau (limite barren plateau), une réduction PCA globale perd de l'information. L'idée
innovante est d'assigner \emph{différentes sous-ensembles} de features à \emph{différents} noyaux
quantiques, puis de combiner les noyaux via MKL.

Le problème d'assignation optimale est NP-difficile (\#vars $= d \times M$, $2^{dM}$ solutions),
mais peut être encodé comme un QUBO et attaqué par QAOA — qui constitue une application candidate
à l'avantage quantique.

\subsection{Pipeline V2}

\begin{figure}[H]
\centering
\begin{tikzpicture}[
  box/.style={rectangle,draw,rounded corners,fill=qgreen!15,minimum width=3.5cm,minimum height=0.9cm,align=center,font=\small},
  arrow/.style={->,thick,qgreen!70!black}
]
\node[box] (fred) at (0,0) {Données FRED\\(19 indicateurs macroéco)};
\node[box] (scale) at (5,0) {Prétraitement\\QuantumScaler};
\node[box] (marg) at (10,0) {Alignements\\marginaux $a_{k,m}$};
\node[box] (qubo) at (10,-2) {Construction QUBO\\(Formulation C)};
\node[box] (solver) at (5,-2) {Solveur QUBO\\(BF / SA / Greedy / QAOA)};
\node[box] (assign) at (0,-2) {Assignation $x_{k,m}$\\features $\to$ noyaux};
\node[box] (mkl2) at (0,-4) {Calcul $M$ noyaux\\sur sous-ensembles};
\node[box] (combine) at (5,-4) {Combinaison MKL\\(Centered Alignment)};
\node[box] (svm2) at (10,-4) {QSVM + Évaluation\\AUC};

\draw[arrow] (fred) -- (scale);
\draw[arrow] (scale) -- (marg);
\draw[arrow] (marg) -- (qubo);
\draw[arrow] (qubo) -- (solver);
\draw[arrow] (solver) -- (assign);
\draw[arrow] (assign) -- (mkl2);
\draw[arrow] (mkl2) -- (combine);
\draw[arrow] (combine) -- (svm2);
\end{tikzpicture}
\caption{Pipeline QMKL-v2 (assignation QUBO)}
\label{fig:v2_pipeline}
\end{figure}

\subsection{Données FRED (Macroéconomiques)}

\texttt{src/data/fred\_loader.py} (200 lignes) charge 19 indicateurs via l'API FRED
(Federal Reserve Economic Data) :

\begin{table}[H]
\centering
\caption{Features FRED — 19 indicateurs macroéconomiques}
\label{tab:fred_features}
\begin{tabular}{@{}llll@{}}
\toprule
\textbf{Groupe} & \textbf{Série FRED} & \textbf{Description} & \textbf{Transform.} \\
\midrule
\multirow{2}{*}{Marché du travail}
  & UNRATE & Taux de chômage & Niveau \\
  & PAYEMS & Emplois non-agricoles & YoY \% \\
\multirow{3}{*}{Monétaire}
  & FEDFUNDS & Taux Fed Funds & Niveau \\
  & DGS10, DGS2 & Rendements 10Y, 2Y & Niveau \\
  & T10Y2Y, T10Y3M & Spreads de taux & bp \\
\multirow{3}{*}{Crédit/Risque}
  & BAMLH0A0HYM2 & Spread High Yield & bp \\
  & BAA10Y & Spread Baa-Treasury & bp \\
  & TEDRATE & Spread TED & bp \\
\multirow{4}{*}{Activité réelle}
  & INDPRO & Production industrielle & YoY \% \\
  & RSAFS & Ventes au détail & YoY \% \\
  & HOUST, PERMIT & Mises en chantier & Niveau \\
\multirow{2}{*}{Inflation/Monnaie}
  & CPIAUCSL & IPC (inflation) & YoY \% \\
  & M2SL, DCOILWTICO & M2, Pétrole WTI & YoY \% \\
\multirow{2}{*}{Sentiment}
  & UMCSENT & Conf. consommateurs U. Mich & Niveau \\
  & SP500 & S\&P 500 & YoY \% \\
\midrule
\multicolumn{2}{l}{\textbf{Cible}} & \multicolumn{2}{l}{USREC (indicateur NBER, 0=expansion, 1=récession)} \\
\bottomrule
\end{tabular}
\end{table}

La série couvre 1970--2023 ($\approx 600$ observations mensuelles). La variable cible USREC
présente un fort déséquilibre ($\approx 15\%$ récessions). Une version synthétique est disponible
hors-ligne via \texttt{load\_fred\_recession\_synthetic()}.

\subsection{Noyaux Analytiques (sans Qiskit)}

Pour permettre des diagnostics rapides, \texttt{src/kernels/analytical.py} (120 lignes) implémente
les formes analytiques fermées \eqref{eq:kz_analytical}--\eqref{eq:kzz_analytical} en NumPy pur :

\begin{lstlisting}[caption=Noyaux analytiques — src/kernels/analytical.py (extrait)]
def kernel_Z(X1, X2, alpha=1.0):
    """Noyau ZFeatureMap analytique. O(N^2 * Q), sans Qiskit."""
    # K[i,j] = prod_k cos^2(alpha * (X1[i,k] - X2[j,k]))
    diff = X1[:, None, :] - X2[None, :, :]   # (N1, N2, Q)
    return np.prod(np.cos(alpha * diff) ** 2, axis=-1)

def kernel_ZZ(X1, X2, alpha=1.0):
    """Noyau ZZFeatureMap analytique. Interactions 2-corps."""
    K_z = kernel_Z(X1, X2, alpha)
    # Termes ZZ: cos^2(alpha * (xi*xj - xi'*xj')) pour k<l
    n1, n2, Q = X1.shape[0], X2.shape[0], X1.shape[1]
    K_zz = np.ones((n1, n2))
    for k in range(Q - 1):
        for l in range(k + 1, Q):
            d_kl = X1[:, None, k] * X1[:, None, l] - X2[None, :, k] * X2[None, :, l]
            K_zz *= np.cos(alpha * d_kl) ** 2
    return K_z * K_zz
\end{lstlisting}

Cette implémentation est $\sim 50\times$ plus rapide que la simulation Qiskit pour les
diagnostics (pas de construction de circuit ni d'overhead simulateur).

\subsection{Implémentation QUBO}

\subsubsection{Calcul des Alignements Marginaux}

\begin{lstlisting}[caption=Alignements marginaux — src/qubo/assignment\_qubo.py (extrait)]
def compute_marginal_alignments(X, y, M, Q, kernel_configs, padding='zero'):
    """
    Retourne a[d, M] : alignement marginal de chaque (feature k, noyau m).
    Pour chaque k: construit X_sub de taille (N, Q) avec feature k + padding,
    calcule K_m(X_sub), puis centered_alignment avec K_target.
    """
    d = X.shape[1]
    a = np.zeros((d, M))
    K_target = (y[:, None] == y[None, :]).astype(float)  # Noyau cible binaire
    K_target_c = center_kernel(K_target)

    for m, cfg in enumerate(kernel_configs):
        kernel_fn = build_kernel_fn(cfg)
        for k in range(d):
            X_sub = build_subset(X, k, Q, padding)
            K_m = kernel_fn(X_sub, X_sub)
            a[k, m] = centered_alignment(K_m, K_target_c)
    return a
\end{lstlisting}

\subsubsection{Construction de la Matrice QUBO}

\begin{lstlisting}[caption=Construction QUBO — src/qubo/assignment\_qubo.py (extrait)]
def build_qubo_matrix(a, d, M, Q, lambda_div=0.5, mu1=2.0, mu2=0.0):
    """
    Construit Q_mat triangulaire sup. taille (d*M, d*M).
    E(x) = x^T @ Q_mat @ x est l'energie QUBO a minimiser.
    """
    n = d * M
    Q_mat = np.zeros((n, n))

    for m in range(M):
        for k in range(d):
            i = k * M + m
            # Diagonal: qualite + terme de contrainte de taille
            Q_mat[i, i] = -a[k, m] + mu1 * (1 - 2 * Q)

            # Meme noyau m, autres features k2 > k: contrainte de taille
            for k2 in range(k + 1, d):
                j = k2 * M + m
                Q_mat[i, j] += 2 * mu1

            # Meme feature k, autres noyaux m2 > m: penalite diversite
            for m2 in range(m + 1, M):
                j = k * M + m2
                Q_mat[i, j] += lambda_div

    return Q_mat
\end{lstlisting}

\subsection{Solveurs QUBO — Implémentation Détaillée}

\subsubsection{Recuit Simulé}

\begin{lstlisting}[caption=Recuit simulé — src/qubo/solvers.py (extrait)]
def solve_simulated_annealing(Q_mat, d, M, Q_size,
                               n_iter=50000, T_init=5.0, T_min=1e-4, seed=42):
    n = d * M
    rng = np.random.RandomState(seed)

    # Initialisation depuis solution faisable (non-chevauchante)
    x = np.zeros(n, dtype=int)
    for m in range(M):
        for q_idx in range(Q_size):
            k = (m * Q_size + q_idx) % d
            x[k * M + m] = 1

    best_x, best_E = x.copy(), energy(x, Q_mat)
    T = T_init
    cooling = (T_min / T_init) ** (1.0 / n_iter)

    for _ in range(n_iter):
        idx = rng.randint(n)
        x[idx] ^= 1                             # Flip bit
        dE = energy(x, Q_mat) - energy(best_x, Q_mat)  # Delta energie

        if dE < 0 or rng.rand() < np.exp(-dE / T):
            if energy(x, Q_mat) < best_E:
                best_x, best_E = x.copy(), energy(x, Q_mat)
        else:
            x[idx] ^= 1                         # Rejeter le mouvement
        T *= cooling

    return {'x': best_x, 'energy': best_E, 'solver': 'sa'}
\end{lstlisting}

\subsubsection{QAOA — Implémentation Locale}

\begin{lstlisting}[caption=QAOA simulation — src/qubo/solvers.py (extrait, version locale)]
def solve_qaoa(Q_mat, d, M, Q_size, p=1, max_iter=50, seed=42, backend=None):
    n_vars = d * M

    # Construction de l'Hamiltonien d'Ising: H_C = sum_{i<=j} Q[i,j] * (I-Zi)/2 * (I-Zj)/2
    pauli_list = qubo_to_ising_hamiltonian(Q_mat)   # -> SparsePauliOp
    hamiltonian = SparsePauliOp.from_list(pauli_list)

    # Circuit QAOA a p couches
    qc = build_qaoa_circuit(n_vars, hamiltonian, p)

    # Optimisation COBYLA
    def cost_fn(theta):
        bound_qc = qc.assign_parameters(theta)
        if backend is None:
            ev = StatevectorEstimator().run([(bound_qc, hamiltonian)]).result()[0].data.evs
        else:
            # Mode hardware: job IBM Quantum (sans Session pour Open Plan)
            estimator_hw = RuntimeEstimatorV2(backend)
            pub = (isa_qc, isa_hamiltonian, [theta])
            job = estimator_hw.run([pub])
            # Sauvegarde de l'ID avant .result() (protection contre timeout)
            save_job_id(job.job_id(), theta)
            ev = job.result()[0].data.evs
        return float(ev)

    # Multi-redemarrages (3 pour p<=2, 1 pour hardware)
    n_restarts = 1 if backend is not None else (3 if p <= 2 else 1)
    best_result = None
    for restart in range(n_restarts):
        theta_init = np.random.RandomState(seed + restart * 13).uniform(0, np.pi/2, 2*p)
        res = minimize(cost_fn, theta_init, method='COBYLA', options={'maxiter': max_iter})
        if best_result is None or res.fun < best_result.fun:
            best_result = res

    # Extraction de la solution: iteration sur tous les bitstrings, selection par energie QUBO
    sampler_result = sample_qaoa_circuit(qc, best_result.x)
    best_x = min(sampler_result.keys(), key=lambda x: energy(x, Q_mat))
    return {'x': best_x, 'energy': energy(best_x, Q_mat), 'solver': 'qaoa'}
\end{lstlisting}

\subsection{Expérience IBM Quantum Hardware}
\label{sec:hardware}

\subsubsection{Configuration}

\begin{itemize}
  \item \textbf{Backend} : IBM Torino (Heron r2, 133 qubits)
  \item \textbf{Plan} : IBM Open Plan (gratuit, sans Sessions)
  \item \textbf{Instance} : $d=12$, $M=3$, $Q=4$ $\to$ 36 variables QUBO $\to$ 36-qubit QAOA
  \item \textbf{Algorithme} : QAOA $p=1$, optimisation COBYLA, 50 itérations demandées
\end{itemize}

\subsubsection{Erreurs Rencontrées et Solutions}

\begin{warningbox}[Obstacles techniques hardware — chronologie des erreurs]

\paragraph{Erreur 1 : Session non disponible sur Open Plan.}
\begin{verbatim}
RequestsApiError: 400 — You are not authorized to run a session
when using the open plan.
\end{verbatim}
\textbf{Solution} : Remplacer \texttt{Session(backend=backend)} par appels directs
\texttt{RuntimeEstimatorV2(backend)} (mode "job unique").

\paragraph{Erreur 2 : Inadéquation qubits circuit/observable.}
\begin{verbatim}
ValueError: The number of qubits of the circuit (133) does not
match the number of qubits of the ()-th observable (36).
\end{verbatim}
\textbf{Cause} : Après transpilation vers le backend 133-qubit, le circuit occupe 133 qubits
physiques, mais l'Hamiltonien n'en a que 36. \\
\textbf{Solution} :
\begin{lstlisting}[numbers=none]
pm = generate_preset_pass_manager(optimization_level=1, backend=backend)
isa_qc_template = pm.run(qc)
isa_hamiltonian = hamiltonian.apply_layout(
    isa_qc_template.layout,
    num_qubits=isa_qc_template.num_qubits,   # 133
)
\end{lstlisting}

\paragraph{Erreur 3 : Dépassement du quota temps Open Plan.}
50 itérations $\times$ 12 s/job = 10 min $>$ 4 min disponibles. \\
\textbf{Solution} : Arrêt forcé du kernel, sauvegarde des IDs de jobs déjà exécutés
dans \texttt{results/qubo\_solutions/hw\_jobs.json}, récupération post-hoc des résultats
via \texttt{service.jobs()}.

\paragraph{Erreur 4 : Simulation statevector 36 qubits.}
$2^{36} = 69 \times 10^9$ états $\times$ 16 octets $= 1$ To RAM $\to$ mort du kernel. \\
\textbf{Solution} : Interdire toute simulation locale pour $n > 20$ variables.
Notebook NB03 averti explicitement.

\end{warningbox}

\subsubsection{Données Récupérées}

24 jobs Estimator ont été exécutés avant l'épuisement du quota. Chaque job calcule
$\langle H_C \rangle_{\bm{\theta}}$ pour une valeur de $\bm{\theta} = (\beta, \gamma)$
soumise par COBYLA.

\begin{table}[H]
\centering
\caption{Données hardware disponibles (jobs Estimator IBM Torino)}
\label{tab:hardware_jobs}
\begin{tabular}{@{}ccccc@{}}
\toprule
\textbf{Type Job} & \textbf{Nombre} & \textbf{Observable} & \textbf{Résultat} & \textbf{Fichier} \\
\midrule
Estimator (EV) & 24 & $H_C$ (Ising, 36 qubits) & $\langle H_C\rangle$ pour 24 $\bm{\theta}$ & hw\_jobs.json \\
Sampler & 0 & — & Bitstrings non récupérés & — \\
\bottomrule
\end{tabular}
\end{table}

\begin{warningbox}[Limitation majeure des résultats hardware]
Sans Sampler (qui aurait nécessité $\sim 10$ jobs supplémentaires inaccessibles), il est
\textbf{impossible d'extraire un bitstring solution}. Les 24 valeurs propres servent
uniquement à tracer la \emph{courbe de convergence COBYLA}, montrant que l'énergie
$\langle H_C \rangle$ \emph{diminue} au fil des itérations — preuve que l'optimiseur
progresse, mais pas la preuve d'une solution finale.

L'assignation finale a été obtenue par Recuit Simulé (SA) — les chiffres AUC du hardware
sont identiques à ceux du SA, ce qui constitue un résultat négatif \emph{sur la
comparabilité}, non sur l'approche elle-même.
\end{warningbox}

\subsubsection{Analyse de la Limite de Simulation}

Le coût de simulation QAOA pour $n$ qubits est :
\begin{equation}
  \text{Coût}_{\text{classique}} = O\!\left(N^2 \cdot 2^n\right) \times n_{\text{iter}}
\end{equation}
Pour $n = 36$ : $2^{36} = 69\,438\,976\,512$ états $\times$ 16 octets $\approx$ 1 To.

La limite pratique est $n \leq 20$--24 qubits sur laptop (16 Go RAM). Au-delà, seul
le matériel quantique (ou un cluster HPC) permet l'évaluation.

\subsection{Notebooks — Architecture de Recherche}

\subsubsection{Trois Notebooks Actifs}

\paragraph{NB01 — Barren Plateaus.}
Démontre empiriquement la concentration exponentielle des noyaux $K_Z$ et $K_{ZZ}$
en fonction de $Q \in \{2, 4, 6, 8, 10, 12\}$ :
\begin{itemize}
  \item Décroissance de $\sigma_{\text{off}}(K)$ (std hors-diagonale)
  \item Décroissance de GradVar($K$, $\alpha$) (variance du gradient)
  \item Impact AUC sur German Credit : dégradation dès $Q \geq 8$
  \item QMKL avec $M=5$ noyaux $\times$ $Q=4$ : récupération de l'AUC
\end{itemize}

\paragraph{NB02 — Complexité Crossover.}
Analyse théorique + empirique du croisement simulation classique / calcul quantique :
\begin{itemize}
  \item Timing simulation classique : $t_{\text{classique}}(Q) \propto 2^Q$ (mesuré)
  \item Comptage ressources quantiques : \#CNOT $\propto Q^2/2$, profondeur $\propto Q$
  \item Projection : la simulation devient plus coûteuse que l'exécution quantique vers
        $Q^* \approx 50$ (hypothèse : hardware futur sans erreur)
\end{itemize}

\paragraph{NB03 — Hardware Torino.}
\begin{itemize}
  \item Chargement de \texttt{hw\_jobs.json} (1 job sauvegardé)
  \item Récupération des 24 EVs depuis compte IBM Cloud
  \item Courbe de convergence COBYLA
  \item Comparaison AUC : SA vs QAOA-HW (proxy SA, car Sampler non disponible)
  \item Simulation Aer avec modèle de bruit (FakeTorinov2)
\end{itemize}

\subsection{Résultats V2}

\begin{table}[H]
\centering
\caption{Résultats QUBO vs baselines — German Credit ($d=20$, $M=5$, $Q=4$)}
\label{tab:v2_results}
\begin{tabular}{@{}lcccc@{}}
\toprule
\textbf{Méthode d'assignation} & \textbf{AUC} & \textbf{$\pm$ Std} & \textbf{Solveur} & \textbf{Temps (s)} \\
\midrule
Single kernel (PCA, $Q=4$)      & $\approx 0.69$ & 0.08 & — & — \\
Random subsets ($Q=4$, $M=5$)   & $0.72$ & 0.05 & Random & — \\
Non-overlap fixe                 & $0.75$ & 0.04 & Déterministe & — \\
QUBO-SA                          & $0.76$ & 0.04 & SA & $\sim 8$ s \\
\textbf{QUBO-QAOA (simulé)}     & \textbf{$0.76$} & 0.04 & QAOA $p=1$ & $\sim 5$ min \\
QMKL v1 (PCA $Q=4$)            & $0.74$ & 0.06 & Centered Align. & — \\
\textbf{RBF-SVM}                & \textbf{$0.80$} & 0.04 & — & — \\
\bottomrule
\end{tabular}
\end{table}

\begin{resultbox}[Résultats cibles vs résultats obtenus]
Les résultats cibles du CLAUDE.md visaient AUC = 0.78--0.82 pour QUBO-assignation.
En pratique, nous obtenons $\approx 0.76$, soit une amélioration réelle par rapport à
Random subsets ($+4$ points) et Single kernel PCA ($+7$ points), mais toujours inférieur
au RBF-SVM classique ($-4$ points).
\end{resultbox}

\subsection{Ce Qui N'a Pas Fonctionné en V2}

\begin{warningbox}[Résultats négatifs V2 — Enseignements]
\begin{enumerate}
  \item \textbf{QAOA simulation 36 qubits} : Impossible sur laptop (1 To RAM).
        La démonstration complète nécessite un HPC ou hardware quantique.

  \item \textbf{IBM Open Plan — Sessions} : L'absence de Sessions force un mode "job unique"
        par itération COBYLA, ajoutant $\sim 10$--12 s de latence réseau par évaluation.
        Avec 50 itérations et un quota 4 min, le test complet est infaisable en Open Plan.

  \item \textbf{Récupération bitstring} : Les 24 jobs Estimator ne donnent que
        $\langle H_C \rangle$, pas de bitstring. Pour extraire la solution binaire,
        un job Sampler supplémentaire est requis — non disponible ici.

  \item \textbf{Données FRED} : La série temporelle avec dépendance temporelle viole
        l'hypothèse i.i.d. des SVM. Un splitting temporel strict est nécessaire
        (pas de validation croisée aléatoire), ce qui réduit effectivement la quantité
        de données d'entraînement.

  \item \textbf{QAOA $p=1$ insuffisant} : Pour 36 variables, un ansatz QAOA de
        profondeur $p=1$ ($\bm{\theta} \in \mathbb{R}^2$) est peu expressif.
        La convergence en énergie est visible mais la solution peut être loin de l'optimum.
\end{enumerate}
\end{warningbox}

% ================================================================
\section{Analyse Comparative et Discussion}
\label{sec:discussion}

\subsection{Quand l'Avantage Quantique se Manifeste-t-il ?}

La question centrale de v2 est empiriquement difficile à répondre aujourd'hui.
Voici une analyse des régimes :

\begin{table}[H]
\centering
\caption{Analyse des régimes d'avantage quantique potentiel}
\label{tab:quantum_advantage}
\begin{tabular}{@{}lcccl@{}}
\toprule
\textbf{Régime} & $Q$ & \textbf{Simulation faisable} & \textbf{Matériel disponible} & \textbf{Statut} \\
\midrule
Démonstration QML & 4--6 & Oui (laptop) & Open Plan & Étudié, pas d'avantage \\
Limite barren plateau & 8--15 & Oui (HPC) & Open/Standard & Dégradation observée \\
Simulation difficile & 20--30 & Non (laptop) & Standard Plan & Non testé \\
Infaisable classiquement & $\geq 50$ & Jamais & Future hardware & Avantage théorique \\
\bottomrule
\end{tabular}
\end{table}

\textbf{Conclusion} : Dans le régime expérimental accessible ($Q \leq 6$, simulation CPU),
QMKL n'offre pas d'avantage par rapport au RBF-SVM. L'avantage potentiel requiert
$Q \geq 50$ sur hardware sans erreur — un horizon à $\sim 5$--10 ans.

\subsection{Comparaison V1 vs V2}

\begin{table}[H]
\centering
\caption{Bilan V1 vs V2}
\label{tab:v1_v2_comparison}
\begin{tabular}{@{}lll@{}}
\toprule
\textbf{Dimension} & \textbf{V1 (Finance)} & \textbf{V2 (QUBO)} \\
\midrule
Innovation principale & Diversité de 12 noyaux + 5 stratégies MKL & Assignation QUBO optimale \\
Dataset principal & German Credit (1000 éch.) & FRED Recession (600 obs.) \\
Stratégie de réduction & PCA globale $\to Q$ dim. & Sous-ensembles $Q$ features/noyau \\
Avantage vs Random & Non significatif & $+4$ points AUC (p=0.08) \\
Avantage vs RBF-SVM & $-7$ à $-12$ points & $-4$ points \\
Test hardware & Non & IBM Torino, 24 jobs \\
Limite principale & Barren plateau ($Q > 6$) & Quota temps IBM Open Plan \\
Contribution réelle & Étude empirique complète & Preuve de concept QUBO \\
\bottomrule
\end{tabular}
\end{table}

\subsection{Architectures Commune des Deux Projets}

Les modules \texttt{src/preprocessing/scaler.py}, \texttt{src/mkl/alignment.py},
\texttt{src/mkl/combiner.py} et \texttt{src/evaluation/metrics.py} sont identiques
dans les deux projets. Cette convergence architecturale confirme la pertinence du
découpage modulaire. En v2, l'ajout de \texttt{src/kernels/analytical.py} (noyaux
NumPy pur) et \texttt{src/qubo/} (module QUBO) étend le pipeline sans rupture.

\subsection{Valeur des Résultats Négatifs}

Les projets QMKL constituent, au-delà des AUC, une documentation précieuse de :
\begin{enumerate}
  \item \textbf{Les barren plateaus en pratique} : Mesures empiriques de $\sigma_{\text{off}}$
        et GradVar confirment la décroissance $\sim 2^{-Q}$.
  \item \textbf{Les obstacles IBM Quantum Open Plan} : Sessions indisponibles, quota temps
        court, latence 12 s/job — barrières concrètes à l'utilisation académique.
  \item \textbf{La complexité de la simulation} : 36 qubits = 1 To RAM — frontière tangible
        entre simulation et hardware.
  \item \textbf{L'implémentation QUBO complète} : De la formulation à la transpilation
        hardware, en passant par le mapping Ising et la correction du layout.
\end{enumerate}

% ================================================================
\section{Infrastructure Technique}
\label{sec:infra}

\subsection{Dépendances}

\begin{table}[H]
\centering
\caption{Stack technologique}
\label{tab:dependencies}
\begin{tabular}{@{}llll@{}}
\toprule
\textbf{Catégorie} & \textbf{Package} & \textbf{Version} & \textbf{Usage} \\
\midrule
\multirow{4}{*}{Quantum}
  & qiskit & $\geq 1.0$ & Circuits, transpilation \\
  & qiskit-aer & $\geq 0.14$ & Simulation CPU/GPU \\
  & qiskit-machine-learning & $\geq 0.7$ & FidelityQuantumKernel, PQK \\
  & qiskit-ibm-runtime & $\geq 0.20$ & IBM Quantum backend, Estimator/Sampler \\
\midrule
\multirow{2}{*}{ML}
  & scikit-learn & $\geq 1.3$ & SVM, PCA, métriques \\
  & scikit-optimize & $\geq 0.9$ & Bayesian Optimization \\
\midrule
Optimisation & cvxpy & $\geq 1.4$ & SDP, QP (MKL weights) \\
\midrule
\multirow{3}{*}{Données}
  & numpy, scipy & — & Algèbre linéaire, stats \\
  & pandas & $\geq 2.0$ & Manipulation des séries FRED \\
  & requests & — & API FRED \\
\midrule
\multirow{2}{*}{Viz.}
  & matplotlib & $\geq 3.7$ & Figures publication \\
  & seaborn & — & Heatmaps, distributions \\
\bottomrule
\end{tabular}
\end{table}

\subsection{Conventions Matplotlib}

\begin{lstlisting}[caption=Conventions graphiques — standard du projet]
plt.rcParams.update({
    'font.family': 'sans-serif',
    'figure.dpi': 150,
    'savefig.bbox': 'tight'
})

COLORS = {
    'Z': '#3498db',     # Bleu
    'ZZ': '#e74c3c',    # Rouge
    'XZ': '#2ecc71',    # Vert
    'YXX': '#f39c12',   # Orange
    'YZX': '#9b59b6',   # Violet
    'Pauli': '#1abc9c', # Turquoise
}
\end{lstlisting}

\subsection{Reproductibilité}

\begin{itemize}
  \item Toutes les graines aléatoires sont fixées (\texttt{seed=42} par défaut)
  \item Les matrices noyaux sont cachées sur disque (MD5 hash des paramètres)
  \item Les résultats hardware sont sauvegardés dans \texttt{results/qubo\_solutions/}
  \item La version synthétique de FRED est disponible hors-ligne
\end{itemize}

% ================================================================
\section{Conclusion}
\label{sec:conclusion}

\subsection{Bilan}

Ce travail a implémenté et évalué deux approches de Quantum Multiple Kernel Learning pour
la classification financière :

\paragraph{V1 (QMKL-Finance).}
Une étude empirique exhaustive sur 12 noyaux, 5 stratégies MKL et 3 datasets montre que
QMKL n'améliore pas le RBF-SVM dans le régime actuel ($Q \leq 6$, simulation classique).
Le phénomène de barren plateau en est la cause principale.

\paragraph{V2 (QMKL-v2).}
L'assignation QUBO de features aux noyaux est une approche innovante qui améliore
$+4$ points AUC par rapport aux sous-ensembles aléatoires sur German Credit.
L'implémentation complète du pipeline QAOA jusqu'au hardware IBM Torino est fonctionnelle,
mais limitée par le quota Open Plan et la taille du problème.

\subsection{Contributions}

\begin{enumerate}
  \item \textbf{Implémentation complète} : 40+ modules Python, 19 notebooks, architecture
        modulaire et réutilisable.
  \item \textbf{Bibliothèque de 12 noyaux} : Diversité Pauli, paramètres $\alpha$ calibrés,
        cache disque, calcul parallèle.
  \item \textbf{Formes analytiques} : $K_Z$, $K_{ZZ}$ en NumPy pur ($50\times$ plus rapide).
  \item \textbf{QUBO Formulation C} : Encodage complet features $\to$ noyaux, 4 solveurs.
  \item \textbf{QAOA hardware} : Configuration Open Plan (sans Sessions), correction layout
        133-qubit, sauvegarde job IDs, récupération post-hoc.
  \item \textbf{Diagnostics barren plateaus} : Métriques $\sigma_{\text{off}}$ et GradVar.
  \item \textbf{Documentation des obstacles} : Session, quota, qubit mismatch, statevector
        overflow — obstacles concrets documentés pour futurs praticiens.
\end{enumerate}

\subsection{Perspectives}

\begin{itemize}
  \item \textbf{Court terme} : Tester avec IBM Standard Plan (Sessions disponibles,
        quota illimité) pour compléter les 50 itérations COBYLA hardware.
  \item \textbf{Moyen terme} : Évaluer sur instances $Q = 20$--30 avec HPC simulateur
        (cluster GPU, Aer GPU backend).
  \item \textbf{Long terme} : Lorsque le matériel quantique atteint $Q \geq 50$ qubits
        fiables (erreur $< 0.1\%$), l'avantage théorique pourrait se matérialiser —
        notamment pour les séries financières à haute dimension où le barren plateau
        est le principal goulot d'étranglement actuel.
\end{itemize}

% ================================================================
\appendix

\section{Structure des Fichiers}
\label{appendix:files}

\subsection{V1 — Projet-QMKL-Finance}

\begin{verbatim}
Projet-QMKL-Finance/
+-- src/
|   +-- kernels/
|   |   +-- quantum_kernel.py     (180 l.) — builders Qiskit
|   |   +-- feature_maps.py       (188 l.) — 12 feature maps
|   |   \-- kernel_matrix.py      (236 l.) — cache + parallélisme
|   +-- mkl/
|   |   +-- alignment.py          (339 l.) — SDP, Centered, Projection
|   |   +-- combiner.py           (108 l.) — MultipleKernelCombiner
|   |   +-- bayesian_optimizer.py (171 l.) — BO via scikit-optimize
|   |   +-- riemannian_mkl.py           — géométrie différentielle
|   |   \-- shapley_mkl.py              — valeurs de Shapley
|   +-- models/
|   |   +-- qsvm.py               (48 l.)  — wrapper QSVM
|   |   \-- classifier.py               — QMKLClassifier orchestrateur
|   +-- preprocessing/
|   |   +-- scaler.py             (34 l.)  — QuantumScaler [0,2pi]
|   |   \-- feature_reduction.py        — PCA wrapper
|   \-- evaluation/
|       +-- metrics.py            (76 l.)  — AUC, F1, recall
|       +-- visualization.py            — utilitaires graphiques
|       +-- ablation.py                 — ablation par composante
|       \-- statistical_analysis.py     — Wilcoxon, bootstrap CI
+-- data/
|   +-- loaders.py                (130 l.) — German Credit, Bank Mktg, BC
|   \-- quantum_dataset.py              — datasets synthétiques
+-- config/
|   \-- default.yaml              (58 l.)  — configuration expériences
+-- notebooks/                    (19 notebooks)
+-- results/
|   +-- presentation/             (10 figures PNG)
|   \-- kernel_cache/             (~80 matrices .npy)
\-- presentation.tex              (22 slides Beamer)
\end{verbatim}

\subsection{V2 — Projet-QMKL-v2}

\begin{verbatim}
Projet-QMKL-v2/
+-- src/
|   +-- data/
|   |   +-- loaders.py                  — datasets standards
|   |   \-- fred_loader.py        (200 l.) — FRED API + synthétique
|   +-- kernels/
|   |   +-- analytical.py         (120 l.) — K_Z, K_ZZ (NumPy pur)
|   |   +-- subset_kernels.py           — noyaux sur sous-ensembles
|   |   \-- diagnostics.py              — sigma_off, GradVar, timing
|   +-- qubo/
|   |   +-- assignment_qubo.py    (150 l.) — CORE: marginal align, QUBO
|   |   \-- solvers.py            (300 l.) — BF, SA, Greedy, QAOA
|   +-- mkl/                            — identique v1
|   \-- evaluation/                     — identique v1
+-- notebooks/
|   +-- NB01_barren_plateau.ipynb       — barren plateaus empiriques
|   +-- NB02_complexity_crossover.ipynb — crossover sim. classique/QC
|   +-- NB03_hardware_torino.ipynb      — IBM Torino + 24 jobs
|   \-- archive/                        — anciens notebooks 01-06
\-- results/
    +-- figures/                        — 12 figures PNG
    \-- qubo_solutions/
        +-- hw_jobs.json                — IDs jobs IBM
        \-- qaoa_hw_d12_M3_Q4.json     — résultats QAOA hardware
\end{verbatim}

\section{Formules de Référence Rapide}
\label{appendix:formulas}

\begin{align}
  K_Z(\bm{x}, \bm{x}') &= \prod_{k=0}^{Q-1} \cos^2(\alpha(x_k - x'_k)) \tag{A.1} \\[6pt]
  K_{ZZ}(\bm{x}, \bm{x}') &= K_Z(\bm{x}, \bm{x}') \cdot \prod_{k<l} \cos^2(\alpha(x_k x_l - x'_k x'_l)) \tag{A.2} \\[6pt]
  K_{\text{comb}} &= \sum_{m=1}^M w_m K_m, \quad \bm{w} \geq 0, \; \|\bm{w}\|_1 = 1 \tag{A.3} \\[6pt]
  \text{Var}\!\left[\nabla_\theta \mathcal{L}\right] &\sim O(2^{-Q}) \quad \text{(barren plateau)} \tag{A.4} \\[6pt]
  E(\bm{x}) &= \bm{x}^T Q_{\text{QUBO}} \bm{x}, \quad \bm{x} \in \{0,1\}^{dM} \tag{A.5} \\[6pt]
  H_C &= \sum_{i \leq j} Q_{ij} \cdot \frac{I - Z_i}{2} \cdot \frac{I - Z_j}{2} \quad \text{(Ising)} \tag{A.6} \\[6pt]
  |\bm{\theta}\rangle &= \prod_{l=1}^p e^{-i\beta_l H_B} e^{-i\gamma_l H_C} |+\rangle^{\otimes n} \quad \text{(QAOA)} \tag{A.7}
\end{align}

\section{Journal des Obstacles Techniques}
\label{appendix:obstacles}

\begin{longtable}{@{}p{4cm}p{5.5cm}p{5.5cm}@{}}
\toprule
\textbf{Obstacle} & \textbf{Cause} & \textbf{Solution} \\
\midrule
\endfirsthead
\multicolumn{3}{c}{\textit{(suite)} } \\
\toprule
\textbf{Obstacle} & \textbf{Cause} & \textbf{Solution} \\
\midrule
\endhead
SDP instable (v1) & Matrice $M$ mal conditionnée & Cascade CLARABEL$\to$ECOS$\to$SCS$\to$forme fermée \\[4pt]
Concentration noyau ($Q>6$) & Barren plateau, variance $\sim 2^{-Q}$ & QMKL avec $M$ noyaux $\times Q$ petit \\[4pt]
Bank Marketing déséquilibré & 12\% positifs & F1 pondéré, pas AUC seule \\[4pt]
QAOA most-frequent vs best-energy & Mode fréquent $\neq$ énergie minimale & Itérer sur tous les bitstrings, garder min QUBO \\[4pt]
Session IBM Open Plan & Plan gratuit sans Sessions & Job mode direct : \texttt{RuntimeEstimatorV2(backend)} \\[4pt]
Mismatch qubits 133 vs 36 & Layout transpilé 133 qubits & \texttt{hamiltonian.apply\_layout(..., num\_qubits=133)} \\[4pt]
Quota 4 min dépassé & 50 iter $\times$ 12 s = 10 min & Arrêt + récupération post-hoc via \texttt{service.jobs()} \\[4pt]
Statevector 36 qubits = 1 To & $2^{36} \approx 69 \cdot 10^9$ états & Simuler $n \leq 20$, hardware pour $n > 20$ \\[4pt]
Bitstring non récupéré & Sampler non exécuté (quota) & Proxy SA pour AUC ; EV = courbe convergence \\[4pt]
Données FRED temporelles & Dépendance temporelle (non i.i.d.) & Splitting temporel strict \\[4pt]
\bottomrule
\caption{Journal exhaustif des obstacles techniques rencontrés}
\label{tab:obstacles}
\end{longtable}

% ================================================================
% Bibliographie
% ================================================================
\begin{thebibliography}{99}

\bibitem{havlicek2019}
Havlíček, V., Córcoles, A. D., Temme, K., et al. (2019).
Supervised learning with quantum-enhanced feature spaces.
\textit{Nature}, 567(7747), 209–212.

\bibitem{schuld2021}
Schuld, M., \& Killoran, N. (2019).
Quantum Machine Learning in Feature Hilbert Spaces.
\textit{Physical Review Letters}, 122(4), 040504.

\bibitem{cortes2012}
Cortes, C., Mohri, M., \& Rostamizadeh, A. (2012).
Algorithms for Learning Kernels Based on Centered Alignment.
\textit{Journal of Machine Learning Research}, 13, 795–828.

\bibitem{lanckriet2004}
Lanckriet, G. R. G., Cristianini, N., Bartlett, P., El Ghaoui, L., \& Jordan, M. I. (2004).
Learning the Kernel Matrix with Semidefinite Programming.
\textit{Journal of Machine Learning Research}, 5, 27–72.

\bibitem{mcclean2018}
McClean, J. R., Boixo, S., Smelyanskiy, V. N., Babbush, R., \& Neven, H. (2018).
Barren Plateaus in Quantum Neural Network Training Landscapes.
\textit{Nature Communications}, 9, 4812.

\bibitem{farhi2014}
Farhi, E., Goldstone, J., \& Gutmann, S. (2014).
A Quantum Approximate Optimization Algorithm.
\textit{arXiv:1411.4028}.

\bibitem{ibm_hsbc_2023}
Pistoia, M., et al. (2023).
Quantum Machine Learning in Financial Classification Tasks.
\textit{arXiv:2207.07677}. IBM / HSBC.

\bibitem{huang2021}
Huang, H.-Y., Broughton, M., Mohseni, M., et al. (2021).
Power of data in quantum machine learning.
\textit{Nature Communications}, 12, 2631.

\bibitem{glick2024}
Glick, J. R., et al. (2024).
Covariant quantum kernels for data with group structure.
\textit{arXiv:2105.03406v3}.

\bibitem{qiskit2024}
Qiskit contributors (2024).
Qiskit: An open-source framework for quantum computing.
\url{https://github.com/Qiskit/qiskit}.

\end{thebibliography}

\end{document}
